% !TeX root = ../../tesi.tex

\section{Studio preliminare e analisi}

\subsection{L’ecosistema delle applicazioni Shopify}
\label{sec:ecosistema-applicazioni}

Descrivo la fase iniziale dello \emph{stage}, dedicata allo studio del funzionamento di
un’applicazione Shopify (autenticazione, \emph{app proxy}, estensioni del tema, differenze tra
Storefront API e Admin API) e degli strumenti per la costruzione di agenti. Distinguo le conoscenze
che avevo già acquisito lavorando sui temi Liquid da quelle che ho dovuto costruire per la prima
volta.

\subsection{Un’analisi dei requisiti iterativa}
\label{sec:requisiti}

Motivo la scelta, condivisa con il \emph{tutor}, di non svolgere un’analisi dei requisiti completa
all’inizio del progetto: il lavoro aveva carattere esplorativo, nessun cliente ne dettava i
requisiti, e una specifica fissata a priori sarebbe diventata presto obsoleta. Descrivo il
procedimento che ho seguito: a partire dagli obiettivi del piano ho abbozzato i casi d’uso
principali e ne ho ricavato i requisiti, che ho rivisto a ogni iterazione insieme al prodotto. Un
diagramma mostra i casi d’uso nella loro forma finale e una tabella riporta i requisiti a cui sono
arrivato, raggruppati per area, ciascuno con un identificativo e l’obiettivo del piano da cui
deriva; ne misuro la copertura nella \secref{sec:quantitativi}.

% DA DEFINIRE: prima i casi d'uso iniziali, poi come raggruppare i requisiti (utente, agente, app...). Solo dopo valutare se partire dalla specifica cancellata il 15/10 in chat-cart.
